\documentclass[11pt]{article}
\usepackage[margin=1in]{geometry}
\usepackage[T1]{fontenc}
\usepackage{lmodern,microtype,amsmath,amssymb,amsthm,hyperref}
\hypersetup{colorlinks=true,linkcolor=blue,citecolor=blue,urlcolor=blue,
pdftitle={An exact period criterion for focal-distance products in elliptic billiards},
pdfauthor={Alec Kriebel}}
\setlength{\emergencystretch}{2em}
\newtheorem{theorem}{Theorem}
\DeclareMathOperator{\sn}{sn}
\DeclareMathOperator{\cn}{cn}
\DeclareMathOperator{\dn}{dn}
\title{An exact period criterion for focal-distance\\products in elliptic billiards}
\author{Alec Kriebel\thanks{Independent researcher.
ORCID: \href{https://orcid.org/0009-0001-9320-500X}{0009-0001-9320-500X}.}}
\date{11 October 2026}
\begin{document}
\maketitle
\begin{abstract}
For a noncircular elliptic billiard with a strictly nested nondegenerate
confocal elliptical caustic, we prove that the product of ordinary distances
from all vertices to either fixed focus is independent of phase if and only
if the least period is \(2\) modulo \(4\).
This proves the least-period interpretation of entry \(k_{114}\) in
Reznik--Garcia--Koiller's list of unproved invariants and classifies every
other primitive period and every repeated traversal.
The proof identifies the complete divisor of the focal-distance factor
in the classical Jacobi parameter. An exact doubled-triangle example
shows why a condition on return length alone is insufficient.
The parametrization, identities and divisor method are established tools.
This AI-assisted research note is unrefereed.
\end{abstract}

\section{The invariant and its exact scope}
Let the boundary and caustic be
\[
 \mathcal E:\frac{x^2}{a^2}+\frac{y^2}{b^2}=1,\qquad
 \mathcal C:\frac{x^2}{A^2}+\frac{y^2}{B^2}=1,
 \quad a>b>0,\quad A>B>0,
\]
with \(a>A\), \(b>B\) and \(a^2-b^2=A^2-B^2=c^2>0\).
Fix a periodic billiard family tangent to \(\mathcal C\), with
\emph{least period} \(n\), vertices \(P_j\), and fixed foci
\(F_\pm=(\pm c,0)\). Write
\[
 D_n^\pm=\prod_{j=0}^{n-1}|P_j-F_\pm|.
\]
These are ordinary positive Euclidean distances.

Reznik, Garcia and Koiller recorded \(D_N^+\) as invariant \(k_{114}\)
for \(N\equiv2\pmod4\), with proof column \emph{?}, in their 2020
preprint and 2021 article \cite[Table~2, p.~345]{source}.
Their introductory setting is a confocal ellipse pair.
The entry does not explicitly settle whether \(N\) is least period or
a return length permitting repetitions. The following theorem decides
both interpretations within that setting.

\begin{theorem}\label{thm:criterion}
For the fixed noncircular confocal pair above, each \(D_n^\pm\) is
independent of phase if and only if \(n\equiv2\pmod4\).
Every admissible primitive star winding and both orientations are included.
For every even \(n\), \(D_n^+=D_n^-\) at each phase.
For an \(r\)-fold traversal, \(r\ge1\), the product is
\((D_n^\pm)^r\), and is constant exactly in the same least-period class.
For a circular boundary the two foci coincide with the center, and
every listed product is \(a^N\).
\end{theorem}

For return length \(N\equiv2\pmod4\), its least period divides \(N\)
and is either odd or \(2\) modulo \(4\). The product is therefore
constant exactly when its least period is even.
Hyperbolic and degenerate caustics are outside the source setting and
this theorem.

\section{The classical real parameter}
Put \(k=c/A\in(0,1)\), \(k'=\sqrt{1-k^2}\), \(K=K(k)\) and
\(K'=K(k')\). All Jacobi functions have \emph{modulus} \(k\).
Stachel's canonical parametrization
\cite[Theorem~4.3 and (4.9), p.~1614]{stachel} is
\begin{equation}\label{eq:param}
 P(u)=(-a\sn u,b\cn u),\quad
 a=A\frac{\dn v}{\cn v},\quad b=\frac{B}{\cn v},\quad
 0<v<K,\quad P_j(u)=P(u+2jv).
\end{equation}
Here \(v\) is uniquely determined by \(\cn v=B/b\); confocality
gives the other axis relation.

For completeness, the geometric bridge can be checked directly.
The tangent at \(Q(U)=(-A\sn U,B\cn U)\) on \(\mathcal C\) is
\[
 -x\sn U/A+y\cn U/B=1.
\]
Jacobi addition \cite[Section~22.8]{dlmf} places both \(P(U-v)\)
and \(P(U+v)\) on it: substitution gives
\[
 \frac{\dn^2v\,\sn^2U+\cn^2U}
 {1-k^2\sn^2U\,\sn^2v}=1.
\]
The contact lies strictly inside \(\mathcal E\), so this is an
actual tangent chord.
The two tangent directions at \(P=(x,y)\) satisfy
\(d^{\mathsf T}Md=0\), with
\[
 M=\begin{pmatrix}B^2-y^2&xy\\xy&A^2-x^2\end{pmatrix}.
\]
The outer normal \(\nu=(x/a^2,y/b^2)\) and its perpendicular
\(t=(-y/b^2,x/a^2)\) obey
\[
 \nu^{\mathsf T}Mt=\frac{xyc^2}{a^2b^2}
 \left(1-x^2/a^2-y^2/b^2\right)=0.
\]
Also \(\det M=A^2B^2(1-x^2/A^2-y^2/B^2)<0\).
Thus the distinct inward tangent rays are interchanged by reflection
across the normal, proving the elastic billiard law.
The real Jacobi amplitude increases strictly through \(2\pi\) over
\(4K\); hence this parameter covers every boundary point and both
oriented caustic-tangent families.

Real injectivity modulo \(4K\) gives the least-period condition
\begin{equation}\label{eq:closure}
 \delta=2v=\frac{4K\tau}{n},\qquad
 \gcd(\tau,n)=1,\qquad 0<\tau<n/2.
\end{equation}
The orientation is positive and \(\tau\) is its turning number.
Neither \(n=1\) nor \(n=2\) can occur with \(0<v<K\).
Elementary ellipse algebra gives
\[
 |P(u)-F_\pm|^2=(a\pm c\sn u)^2.
\]
Since \(a>c\), the linear factors are strictly positive on the real
axis. Therefore the actual distance functions are
\begin{equation}\label{eq:factor}
 g_\pm(u)=a\pm c\sn u,\qquad
 D_n^\pm(u)=\prod_{j=0}^{n-1}g_\pm(u+j\delta).
\end{equation}

\section{The complete divisor and the period criterion}
Regard the focal factor on
\[
 T=\mathbb C/(4K\mathbb Z+2iK'\mathbb Z),\qquad p=iK'.
\]
This torus is used for \(\sn\) and \(g_\pm\); only the real vertex
period \(4K\) was used in \eqref{eq:closure}.
The standard Jacobi periods, poles and shifts
\cite[Tables~22.4.1 and 22.4.3]{dlmf} give two simple poles of
\(g_+\), at \(p,p+2K\), and
\[
 \sn(p+K\pm v)=\frac{\dn v}{k\cn v}=\frac ac,\qquad
 \sn(u+2K)=-\sn u.
\]
Thus its zeros are \(q_\pm=p-K\pm v\).
They are distinct and nonpolar because \(0<v<K\).
The coefficient \(c\ne0\) prevents removal of either pole.
An elliptic function has equal total zero and pole degree, so these
two zeros are simple and exhaustive.

Let \(H=\langle\delta\rangle\) on \(T\). It has exact order \(n\).
Subtracting \(p\) and dividing real offsets by \(4K\) identifies
it with \(\{j/n\bmod1\}\).
Orders in a product add, so the complete product divisor is the sum
of translates of the divisor of \(g_+\).
The two zeros differ by one step \(\delta\); their translates make
a single zero coset of multiplicity two, represented by
\[
 z=-\frac14+\frac{\tau}{2n}.
\]
The pole cosets are \(H\) and \(1/2+H\).

If \(n\) is odd, these pole cosets are disjoint. Both
\[
 nz=\frac{-n+2\tau}{4},\qquad
 n(z-1/2)=\frac{-3n+2\tau}{4}
\]
have odd numerators and are not integers. The zero coset misses
both pole cosets, leaving \(2n\) simple poles.

If \(n\) is even, \(\tau\) is odd and \(1/2\in H\).
The pole cosets coincide, giving \(n\) double poles.
The double-zero coset cancels them precisely when
\[
 nz=\frac{-n+2\tau}{4}\in\mathbb Z,
\]
equivalently \(n\equiv2\pmod4\).
Then the full divisor vanishes, the product extends holomorphically
to the compact connected torus, and it is constant.
For \(n\equiv0\pmod4\), it retains \(n\) double poles.

The products are holomorphic near the real axis. If a product with
uncancelled poles were constant on the real family, the identity
theorem would force its meromorphic extension to be constant,
a contradiction. This proves physical real-family nonconstancy:
\eqref{eq:factor} equals the ordinary Euclidean product itself
on that line.

Finally \(g_-(u)=g_+(u+2K)\). The cyclic product is invariant under
\(H\), proving the same classification for both foci; when \(n\)
is even, \(2K\in H\) also proves equality at each phase.
For \(n=2m\equiv2\pmod4\), antipodal pairing and phase-zero evaluation
give the positive value
\[
 D_n^\pm=\prod_{j=0}^{m-1}
 \bigl(a^2-c^2\sn^2(2jv)\bigr)>0.
\]
An \(r\)-fold traversal raises the positive primitive product to
the power \(r\). This operation is injective on positive reals,
so it preserves constancy in both directions.
For a circle every focal distance is \(a\).
This proves Theorem~\ref{thm:criterion}.

\section{An exact return-length counterexample}
Take \(a=\sqrt{15}\), \(b=\sqrt7\), \(c=\sqrt8\),
\(A^2=135/16\), \(B^2=7/16\).
The triangle
\[
 T=\bigl((a,0),(-3a/4,7/4),(-3a/4,-7/4)\bigr)
\]
and its central reflection are primitive three-bounce members of
one confocal family. Their vertices lie on \(\mathcal E\).
The side lines are \(x=-3a/4=-A\) and \(x\pm ay=a\);
tangency follows from \(A^2+a^2B^2=a^2\).
At the upper off-axis vertex the inward unit rays
\((a/4,-1/4)\), \((0,-1)\) sum to five times the inward normal
\((a/20,-1/4)\). Symmetry handles the other vertices, independently
checking physical reflection.

At \(F_+=(c,0)\) the primitive products are
\[
 (a-c)(a+3c/4)^2=\frac{15a}{2}+3c,\qquad
 (a+c)(a-3c/4)^2=\frac{15a}{2}-3c.
\]
Twice traversing each triangle gives return length six, but the
six-event products differ by
\[
 \left(\frac{15a}{2}+3c\right)^2-
 \left(\frac{15a}{2}-3c\right)^2=90ac=180\sqrt{30}>0.
\]
Thus listed return length cannot replace least period.

\section{Attribution, verification and status}
The theorem proves the least-period \(k_{114}\) observation and
classifies its broader return-period interpretation, which is partly
false. Stachel's parametrization, classical Jacobi identities and
pole-removal method are prior machinery.
Meromorphic Poncelet functions and pole-removal arguments were used
by Schwartz \cite{schwartz} and, for example,
Akopyan--Schwartz--Tabachnikov \cite{ast} and
Roitman--Garcia--Reznik \cite{bicentric}.
The target-specific step is the focal divisor and its coset arithmetic;
no new elliptic-function method is claimed.
Ferudun's related unrefereed note \cite{ferudun} treats the different
outer-tangent-polygon product \(k_{115}\).

The accompanying priority note records the primary sources checked
and search limits. No explicit earlier proof of this exact target was
located in the checked sources; the novelty of this specific formulation
and exclusive historical priority are not certified.
The full text of a related inertia-and-billiards paper by
Dragovi\'c and Gaji\'c \cite{dg} was unavailable; its uninspected
billiard sections remain a disclosed bibliography coverage gap.
Portable exact checks reproduce finite divisor inventories and
physical counterexample arithmetic. They support the written proof,
and do not replace its all-period analytic argument.
No validated numerical certification or proof-assistant formalization
is claimed.

AI tools were used extensively in solving, drafting, comparing
sources, reproducing computations and independent adversarial
reviews. This preprint has not undergone conventional human peer
review or refereeing. The author is responsible for its claims.

\begin{thebibliography}{9}
\bibitem{source} D.~Reznik, R.~Garcia and J.~Koiller,
Fifty New Invariants of \(N\)-Periodics in the Elliptic Billiard,
\emph{Arnold Mathematical Journal} \textbf{7} (2021), 341--355.
\href{https://doi.org/10.1007/s40598-021-00174-y}{doi:10.1007/s40598-021-00174-y}.
Preprint: \href{https://arxiv.org/abs/2004.12497v11}{arXiv:2004.12497v11}.
\bibitem{stachel} H.~Stachel, On the motion of billiards in ellipses,
\emph{European Journal of Mathematics} \textbf{8} (2022), 1602--1622.
\href{https://doi.org/10.1007/s40879-021-00524-2}{doi:10.1007/s40879-021-00524-2}.
\bibitem{dlmf} F.~W.~J. Olver et al., eds.,
\emph{NIST Digital Library of Mathematical Functions},
Chapter~22: Jacobian elliptic functions,
\url{https://dlmf.nist.gov/22}.
\bibitem{ast} A.~Akopyan, R.~Schwartz and S.~Tabachnikov,
Billiards in ellipses revisited,
\emph{European Journal of Mathematics} \textbf{8} (2022), 1313--1327
(published online in 2020).
\href{https://doi.org/10.1007/s40879-020-00426-9}{doi:10.1007/s40879-020-00426-9}.
\bibitem{bicentric} P.~Roitman, R.~Garcia and D.~Reznik,
New Invariants of Poncelet--Jacobi Bicentric Polygons,
\emph{Arnold Mathematical Journal} \textbf{7} (2021), 619--637.
\href{https://doi.org/10.1007/s40598-021-00188-6}{doi:10.1007/s40598-021-00188-6}.
\bibitem{schwartz} R.~Schwartz,
The pentagram integrals for Poncelet families,
\emph{Journal of Geometry and Physics} \textbf{87} (2015), 432--449.
\href{https://doi.org/10.1016/j.geomphys.2014.07.024}{doi:10.1016/j.geomphys.2014.07.024}.
\bibitem{ferudun} A.~Ferudun,
Focal Distance Products of Outer Elliptic Billiard Polygons,
unrefereed preprint, 2 October 2026.
\href{https://doi.org/10.5281/zenodo.23093258}{doi:10.5281/zenodo.23093258}.
\bibitem{dg} V.~Dragovi\'c and B.~Gaji\'c,
Points with rotational ellipsoids of inertia, envelopes of hyperplanes
which equally fit the system of points in \(\mathbb R^k\), and ellipsoidal billiards,
\emph{Physica D} \textbf{451} (2023), 133776.
\href{https://doi.org/10.1016/j.physd.2023.133776}{doi:10.1016/j.physd.2023.133776}.
\end{thebibliography}
\end{document}
